%% Proposed Solution --- revised benefits and observability sections

%% ────────────────────────────────────────────────────────────────────────────
\section{System Advantages}

The proposed solution provides value beyond technical traceability. It improves
how teams understand change impact, coordinate delivery decisions, and maintain
confidence in the information used during development and release activities.

\subsection{Centralized Traceability and Visibility}

The system consolidates dependency information that is normally scattered across
multiple Azure DevOps views into a single analytical entry point. Instead of
manually navigating between work items, pull requests, branches, commits, and
tags, users obtain a coherent representation of how these artifacts relate to
one another.

This centralization improves day-to-day operational visibility and reduces the
risk of overlooking a critical link in the dependency chain.

\subsection{Faster and Better-Informed Decision-Making}

By reconstructing the dependency graph in a structured form, the system supports
decisions that are often difficult to make with fragmented information alone.
Teams can more quickly assess whether a change is isolated or part of a broader
chain, whether a work item is fully represented in source control activity, and
whether a release candidate contains the expected implementation elements.

The result is not only better traceability, but also better decision support for
integration, validation, and release preparation.

\subsection{Scalability for Large Project Portfolios}

The proposed architecture is designed to remain effective when the number of
artifacts grows. Through batch-oriented retrieval and grouped processing, the
system limits repetitive calls to external services and keeps the resolution
process manageable even for large dependency sets.

This makes the solution suitable not only for isolated analyses, but also for
recurrent use in complex enterprise environments where many repositories and
linked work items must be examined together.

\subsection{Consistency and Reliability of Analysis Results}

The normalization of external data into internal DTOs ensures that the system
produces results in a uniform structure, regardless of differences in Azure
DevOps response formats. This consistency improves the reliability of the final
analysis and helps both backend and UI layers operate on stable, predictable
information.

From a user perspective, this means that dependency reports remain easier to
interpret, compare, and reuse across different analysis scenarios.

\subsection{Extensibility and Long-Term Maintainability}

The modular decomposition of the solution allows new dependency dimensions to be
introduced progressively without redesigning the entire system. Additional data
sources such as pipeline executions, build artifacts, release records, or
deployment environments can be integrated as complementary modules around the
existing core.

This architectural flexibility protects the solution from becoming tightly bound
to its initial scope and supports gradual evolution as organizational needs grow.

\subsection{Improved Collaboration Between Stakeholders}

The solution creates a shared analytical view for developers, release managers,
and technical leads. Because all of them rely on the same reconstructed
dependency model, discussions about readiness, impact, and missing links can be
based on a common factual reference rather than on manual interpretation.

This contributes to smoother coordination and reduces ambiguity during critical
project phases.

%% ────────────────────────────────────────────────────────────────────────────
\section{Operational Observability and Logging Strategy}

Because the proposed solution reconstructs dependency chains from several linked
artifacts, operational transparency is essential. For that reason, the system
includes a structured logging strategy based on \textbf{AppLogger}, which is used
as the standard application logging mechanism across the backend services.

At the proposed-solution level, logging is not treated as a low-level technical
detail, but as a reliability mechanism that helps teams understand how a request
was processed, where a resolution path failed, and which step of the dependency
pipeline produced a given result.

\subsection{Role of Logging in the Solution}

The logging strategy supports four practical objectives:

\begin{itemize}
  \item improving operational transparency during dependency analysis;
  \item facilitating diagnosis when external Azure DevOps calls fail or return incomplete data;
  \item making execution flows easier to follow across the main components;
  \item strengthening confidence in the produced analytical output.
\end{itemize}

\subsection{Log Level Strategy}

To keep the system observable without overwhelming users or operators with
excessive detail, the solution relies on differentiated log levels.

\begin{itemize}
  \item \textbf{log}: records the main execution steps and important state transitions, such as the start and completion of major analysis operations;
  \item \textbf{warn}: highlights recoverable anomalies, partial data situations, or unexpected states that do not fully block processing;
  \item \textbf{error}: captures failures that prevent a dependency resolution step from being completed correctly;
  \item \textbf{debug}: provides additional diagnostic detail useful during investigation and controlled troubleshooting sessions;
  \item \textbf{verbose}: offers very fine-grained tracing reserved for advanced analysis or temporary deep inspection.
\end{itemize}

This level-based approach ensures that the system can remain readable in normal
operation while still offering deeper visibility when investigation is required.

\subsection{Architectural Value of the Logging Strategy}

Within the proposed solution, structured logging plays an important support role
for maintainability and trust. It helps confirm that the dependency pipeline is
executing as expected, highlights where external dependencies affect the result,
and provides a traceable operational narrative for each analysis request.

In this sense, AppLogger contributes not only to technical monitoring, but also
to the overall robustness of the solution by making the behavior of the system
more understandable, auditable, and controllable over time.